
\documentclass{article}
\usepackage{colortbl}
\usepackage{makecell}
\usepackage{multirow}
\usepackage{supertabular}

\begin{document}

\newcounter{utterance}

\centering \large Interaction Transcript for game `dond', experiment `semi\_en', episode 4 with Qwen3.5{-}9B.
\vspace{24pt}

{ \footnotesize  \setcounter{utterance}{1}
\setlength{\tabcolsep}{0pt}
\begin{supertabular}{c@{$\;$}|p{.15\linewidth}@{}p{.15\linewidth}p{.15\linewidth}p{.15\linewidth}p{.15\linewidth}p{.15\linewidth}}
    \# & $\;$A & \multicolumn{4}{c}{Game Master} & $\;\:$B\\
    \hline

    \theutterance \stepcounter{utterance}  
    & & \multicolumn{4}{p{0.6\linewidth}}{
        \cellcolor[rgb]{0.9,0.9,0.9}{
            \makecell[{{p{\linewidth}}}]{
                \texttt{\tiny{[P1$\langle$GM]}}
                \texttt{You are playing a negotiation game in which you have to agree on how to divide a set of items among you and another player.} \\
\\ 
\texttt{Rules:} \\
\texttt{(a) You and the other player are given a set of items. Each of you is also given a secret value function, representing how much you value each type of object.} \\
\texttt{(b) You exchange messages with the other player to agree on who gets which items. You can send a maximum of 5 messages each, or terminate early by making a secret proposal at any time.} \\
\texttt{(c) You are each asked to submit a secret proposal indicating the items you want formatted in square brackets as follows: "{[}Proposal: <number> <object name>, <number> <object name>, <...>{]}"} \\
\texttt{(d) If your proposals are complementary, i.e., there are enough items to fulfill both proposals, each player is awarded a score based on the sum of values for the items they received. Otherwise, both of you get zero points.} \\
\texttt{(e) Your goal is to maximize the score you receive.} \\
\\ 
\texttt{Let us start.} \\
\\ 
\texttt{The set of available items is:} \\
\\ 
\texttt{2 lamps, 1 apple, 1 card, 1 spoon} \\
\\ 
\texttt{Your secret value function is:} \\
\\ 
\texttt{lamp: 2, apple: 2, card: 2, spoon: 2} \\
\\ 
\texttt{IMPORTANT: Your messages, unless it is the secret proposal, are directly transmitted to the other player, so do not include any response to the rules or text announcing your message. To make a secret proposal, use the indicated format. Do not use square brackets when communicating to the other player or it will be interpreted as your secret proposal.} \\
            }
        }
    }
    & \\ \\

    \theutterance \stepcounter{utterance}  
    & \multicolumn{4}{p{0.6\linewidth}}{
        \cellcolor[rgb]{0.9,0.9,0.9}{
            \makecell[{{p{\linewidth}}}]{
                \texttt{\tiny{[P1$\rangle$GM]}}
                \texttt{I value all items equally at 2 points each. The total value of the set is \$2 \textbackslash{}times 2 + 1 \textbackslash{}times 2 + 1 \textbackslash{}times 2 + 1 \textbackslash{}times 2 = 10\$ points. To maximize my score, I need to secure the items with the highest total value. Since I value everything equally, I should aim for a split that gives me the most items or the most valuable combination.} \\
\\ 
\texttt{The other player also has a secret value function, but since I don't know it yet, I need to propose a division that is likely to be complementary. A fair and efficient split would be to give the other player half of the items and keep the other half.} \\
\\ 
\texttt{Let's propose a split where I take the two lamps and the apple, and the other player takes the card and the spoon.} \\
\texttt{My items: 2 lamps, 1 apple.} \\
\texttt{Other player's items: 1 card, 1 spoon.} \\
\\ 
\texttt{This uses 4 items total, leaving none unused. If the other player accepts this (by proposing the complementary set), I get \$2+2+2=6\$ points. If I try to take everything, we likely won't be complementary, and I get 0.} \\
\\ 
\texttt{I will send a message proposing this specific division.} \\
\\ 
\texttt{Proposal: I want the 2 lamps and the 1 apple. The other player should get the 1 card and the 1 spoon.} \\
\\ 
\texttt{Message to other player:} \\
\texttt{I propose that I take the 2 lamps and the 1 apple, and you take the 1 card and the 1 spoon.} \\
            }
        }
    }
    & & \\ \\

    \theutterance \stepcounter{utterance}  
    & & & \multicolumn{2}{p{0.3\linewidth}}{
        \cellcolor[rgb]{0.9,0.9,0.9}{
            \makecell[{{p{\linewidth}}}]{
                \texttt{\tiny{[GM$|$GM]}}
                \texttt{continue} \\
            }
        }
    }
    & & \\ \\

    \theutterance \stepcounter{utterance}  
    & & & \multicolumn{2}{p{0.3\linewidth}}{
        \cellcolor[rgb]{0.9,0.9,0.9}{
            \makecell[{{p{\linewidth}}}]{
                \texttt{\tiny{[GM$|$GM]}}
                \texttt{I value all items equally at 2 points each. The total value of the set is \$2 \textbackslash{}times 2 + 1 \textbackslash{}times 2 + 1 \textbackslash{}times 2 + 1 \textbackslash{}times 2 = 10\$ points. To maximize my score, I need to secure the items with the highest total value. Since I value everything equally, I should aim for a split that gives me the most items or the most valuable combination.} \\
\\ 
\texttt{The other player also has a secret value function, but since I don't know it yet, I need to propose a division that is likely to be complementary. A fair and efficient split would be to give the other player half of the items and keep the other half.} \\
\\ 
\texttt{Let's propose a split where I take the two lamps and the apple, and the other player takes the card and the spoon.} \\
\texttt{My items: 2 lamps, 1 apple.} \\
\texttt{Other player's items: 1 card, 1 spoon.} \\
\\ 
\texttt{This uses 4 items total, leaving none unused. If the other player accepts this (by proposing the complementary set), I get \$2+2+2=6\$ points. If I try to take everything, we likely won't be complementary, and I get 0.} \\
\\ 
\texttt{I will send a message proposing this specific division.} \\
\\ 
\texttt{Proposal: I want the 2 lamps and the 1 apple. The other player should get the 1 card and the 1 spoon.} \\
\\ 
\texttt{Message to other player:} \\
\texttt{I propose that I take the 2 lamps and the 1 apple, and you take the 1 card and the 1 spoon.} \\
            }
        }
    }
    & & \\ \\

    \theutterance \stepcounter{utterance}  
    & & \multicolumn{4}{p{0.6\linewidth}}{
        \cellcolor[rgb]{0.9,0.9,0.9}{
            \makecell[{{p{\linewidth}}}]{
                \texttt{\tiny{[GM$\rangle$P2]}}
                \texttt{You are playing a negotiation game in which you have to agree on how to divide a set of items among you and another player.} \\
\\ 
\texttt{Rules:} \\
\texttt{(a) You and the other player are given a set of items. Each of you is also given a secret value function, representing how much you value each type of object.} \\
\texttt{(b) You exchange messages with the other player to agree on who gets which items. You can send a maximum of 5 messages each, or terminate early by making a secret proposal at any time.} \\
\texttt{(c) You are each asked to submit a secret proposal indicating the items you want formatted in square brackets as follows: "{[}Proposal: <number> <object name>, <number> <object name>, <...>{]}"} \\
\texttt{(d) If your proposals are complementary, i.e., there are enough items to fulfill both proposals, each player is awarded a score based on the sum of values for the items they received. Otherwise, both of you get zero points.} \\
\texttt{(e) Your goal is to maximize the score you receive.} \\
\\ 
\texttt{Let us start.} \\
\\ 
\texttt{The set of available items is:} \\
\\ 
\texttt{2 lamps, 1 apple, 1 card, 1 spoon} \\
\\ 
\texttt{Your secret value function is:} \\
\\ 
\texttt{lamp: 0, apple: 2, card: 3, spoon: 5} \\
\\ 
\texttt{IMPORTANT: Your messages, unless it is the secret proposal, are directly transmitted to the other player, so do not include any response to the rules or text announcing your message. To make a secret proposal, use the indicated format. Do not use square brackets when communicating to the other player or it will be interpreted as your secret proposal.} \\
\\ 
\\ 
\texttt{I value all items equally at 2 points each. The total value of the set is \$2 \textbackslash{}times 2 + 1 \textbackslash{}times 2 + 1 \textbackslash{}times 2 + 1 \textbackslash{}times 2 = 10\$ points. To maximize my score, I need to secure the items with the highest total value. Since I value everything equally, I should aim for a split that gives me the most items or the most valuable combination.} \\
\\ 
\texttt{The other player also has a secret value function, but since I don't know it yet, I need to propose a division that is likely to be complementary. A fair and efficient split would be to give the other player half of the items and keep the other half.} \\
\\ 
\texttt{Let's propose a split where I take the two lamps and the apple, and the other player takes the card and the spoon.} \\
\texttt{My items: 2 lamps, 1 apple.} \\
\texttt{Other player's items: 1 card, 1 spoon.} \\
\\ 
\texttt{This uses 4 items total, leaving none unused. If the other player accepts this (by proposing the complementary set), I get \$2+2+2=6\$ points. If I try to take everything, we likely won't be complementary, and I get 0.} \\
\\ 
\texttt{I will send a message proposing this specific division.} \\
\\ 
\texttt{Proposal: I want the 2 lamps and the 1 apple. The other player should get the 1 card and the 1 spoon.} \\
\\ 
\texttt{Message to other player:} \\
\texttt{I propose that I take the 2 lamps and the 1 apple, and you take the 1 card and the 1 spoon.} \\
            }
        }
    }
    & \\ \\

    \theutterance \stepcounter{utterance}  
    & & & \multicolumn{4}{p{0.6\linewidth}}{
        \cellcolor[rgb]{0.9,0.9,0.9}{
            \makecell[{{p{\linewidth}}}]{
                \texttt{\tiny{[GM$\langle$P2]}}
                \texttt{I propose that I take the 1 card and the 1 spoon, and you take the 2 lamps and the 1 apple.} \\
            }
        }
    }
     \\ \\

    \theutterance \stepcounter{utterance}  
    & & & \multicolumn{2}{p{0.3\linewidth}}{
        \cellcolor[rgb]{0.9,0.9,0.9}{
            \makecell[{{p{\linewidth}}}]{
                \texttt{\tiny{[GM$|$GM]}}
                \texttt{continue} \\
            }
        }
    }
    & & \\ \\

    \theutterance \stepcounter{utterance}  
    & & & \multicolumn{2}{p{0.3\linewidth}}{
        \cellcolor[rgb]{0.9,0.9,0.9}{
            \makecell[{{p{\linewidth}}}]{
                \texttt{\tiny{[GM$|$GM]}}
                \texttt{I propose that I take the 1 card and the 1 spoon, and you take the 2 lamps and the 1 apple.} \\
            }
        }
    }
    & & \\ \\

    \theutterance \stepcounter{utterance}  
    & & \multicolumn{4}{p{0.6\linewidth}}{
        \cellcolor[rgb]{0.9,0.9,0.9}{
            \makecell[{{p{\linewidth}}}]{
                \texttt{\tiny{[P1$\langle$GM]}}
                \texttt{I propose that I take the 1 card and the 1 spoon, and you take the 2 lamps and the 1 apple.} \\
            }
        }
    }
    & \\ \\

    \theutterance \stepcounter{utterance}  
    & \multicolumn{4}{p{0.6\linewidth}}{
        \cellcolor[rgb]{0.9,0.9,0.9}{
            \makecell[{{p{\linewidth}}}]{
                \texttt{\tiny{[P1$\rangle$GM]}}
                \texttt{My secret proposal:} \\
\texttt{{[}Proposal: 2 lamp, 1 apple{]}} \\
            }
        }
    }
    & & \\ \\

    \theutterance \stepcounter{utterance}  
    & & & \multicolumn{2}{p{0.3\linewidth}}{
        \cellcolor[rgb]{0.9,0.9,0.9}{
            \makecell[{{p{\linewidth}}}]{
                \texttt{\tiny{[GM$|$GM]}}
                \texttt{proposal} \\
            }
        }
    }
    & & \\ \\

    \theutterance \stepcounter{utterance}  
    & & & \multicolumn{2}{p{0.3\linewidth}}{
        \cellcolor[rgb]{0.9,0.9,0.9}{
            \makecell[{{p{\linewidth}}}]{
                \texttt{\tiny{[GM$|$GM]}}
                [2, 1, 0, 0]
            }
        }
    }
    & & \\ \\

    \theutterance \stepcounter{utterance}  
    & & \multicolumn{4}{p{0.6\linewidth}}{
        \cellcolor[rgb]{0.9,0.9,0.9}{
            \makecell[{{p{\linewidth}}}]{
                \texttt{\tiny{[GM$\rangle$P2]}}
                \texttt{The other player has submitted a secret proposal. It is now your turn to submit a proposal of your own.} \\
            }
        }
    }
    & \\ \\

    \theutterance \stepcounter{utterance}  
    & & & \multicolumn{4}{p{0.6\linewidth}}{
        \cellcolor[rgb]{0.9,0.9,0.9}{
            \makecell[{{p{\linewidth}}}]{
                \texttt{\tiny{[GM$\langle$P2]}}
                \texttt{{[}Proposal: 1 card, 1 spoon{]}} \\
            }
        }
    }
     \\ \\

    \theutterance \stepcounter{utterance}  
    & & & \multicolumn{2}{p{0.3\linewidth}}{
        \cellcolor[rgb]{0.9,0.9,0.9}{
            \makecell[{{p{\linewidth}}}]{
                \texttt{\tiny{[GM$|$GM]}}
                \texttt{proposal} \\
            }
        }
    }
    & & \\ \\

    \theutterance \stepcounter{utterance}  
    & & & \multicolumn{2}{p{0.3\linewidth}}{
        \cellcolor[rgb]{0.9,0.9,0.9}{
            \makecell[{{p{\linewidth}}}]{
                \texttt{\tiny{[GM$|$GM]}}
                [0, 0, 1, 1]
            }
        }
    }
    & & \\ \\

    \theutterance \stepcounter{utterance}  
    & & & \multicolumn{2}{p{0.3\linewidth}}{
        \cellcolor[rgb]{0.9,0.9,0.9}{
            \makecell[{{p{\linewidth}}}]{
                \texttt{\tiny{[GM$|$GM]}}
                [[2, 1, 0, 0], [0, 0, 1, 1]]
            }
        }
    }
    & & \\ \\

\end{supertabular}
}

\end{document}
